\medskip

\begin{minipage}{0.58\linewidth}
Sur la figure ci-contre :

\begin{itemize}[label=$\bullet~~$]
\item ABCD est un carré de côté 9 cm; 
\item les segments de même longueur sont codés.
\end{itemize}

\begin{enumerate}
\item 
	\begin{enumerate}
		\item Le polygone IJKLMNOP est-il régulier, c'est-à-dire a-t-il
tous ses côtés de même longueur ?

\textbf{Justifier la réponse.}
		\item Justifier que l'aire de la surface IJKLMNOP grisée sur
la figure ci-contre est égale à 63~cm$^2$.
	\end{enumerate}
\end{enumerate}
\end{minipage}\hfill
\begin{minipage}{0.38\linewidth}
\psset{unit=0.9cm}
\begin{pspicture}(0,-1)(5.8,5.8)
%\psgrid
\psframe(0.2,0.2)(5.6,5.6)
\pspolygon[fillstyle=solid,fillcolor=lightgray](2,5.6)(3.8,5.6)(5.6,3.8)(5.6,2)(3.8,0.2)(2,0.2)(0.2,2)(0.2,3.8)%IJKLMNOP
\def\dblv{\psline(-0.05,0.1)(-0.05,-0.1)\psline(0.05,0.1)(0.05,-0.1)}
\rput(2.9,5.6){\dblv}\rput(1,5.6){\dblv}\rput(4.7,5.6){\dblv}
\rput(2.9,0.2){\dblv}\rput(1,0.2){\dblv}\rput(4.7,0.2){\dblv}
\def\dblh{\psline(-0.1,0.05)(0.1,0.05)\psline(-0.1,-0.05)(0.1,-0.05)}
\rput(0.2,4.7){\dblh}\rput(0.2,2.9){\dblh}\rput(0.2,1.1){\dblh}
\rput(5.6,4.7){\dblh}\rput(5.6,2.9){\dblh}\rput(5.6,1.1){\dblh}
\uput[ul](0.2,5.6){A} \uput[ur](5.6,5.6){B} \uput[dr](5.6,0.2){C} \uput[dl](0.2,0.2){D}
\uput[u](2,5.6){I} \uput[u](3.8,5.6){J} \uput[r](5.6,3.8){K} \uput[r](5.6,2){L}
\uput[d](3.8,0.2){M} \uput[d](2,0.2){N} \uput[l](0.2,2){O} \uput[l](0.2,3.8){P}
\end{pspicture}
\end{minipage}

\begin{minipage}{0.58\linewidth}
\begin{enumerate}[start=2]
\item Les diagonales du carré ABCD se coupent en S.

\smallskip

On a tracé le cercle de centre S et de diamètre 9 cm.
	\begin{enumerate}
		\item Déterminer l'aire du disque de centre S et de
diamètre 9 cm.
		\item Montrer que la différence entre l'aire du polygone
IJKLMNOP et l'aire du disque représente moins de
1\,\% de l'aire du disque.
	\end{enumerate}
\end{enumerate}
\end{minipage}\hfill
\begin{minipage}{0.38\linewidth}
\psset{unit=0.9cm}
\begin{pspicture}(0,0)(5.8,5.8)
%\psgrid
\psframe(0.2,0.2)(5.6,5.6)
\pspolygon(2,5.6)(3.8,5.6)(5.6,3.8)(5.6,2)(3.8,0.2)(2,0.2)(0.2,2)(0.2,3.8)%IJKLMNOP
\pscircle[fillstyle=solid,fillcolor=lightgray](2.9,2.9){2.7}
\def\dblv{\psline(-0.05,0.1)(-0.05,-0.1)\psline(0.05,0.1)(0.05,-0.1)}
\rput(2.9,5.6){\dblv}\rput(1,5.6){\dblv}\rput(4.7,5.6){\dblv}
\rput(2.9,0.2){\dblv}\rput(1,0.2){\dblv}\rput(4.7,0.2){\dblv}
\def\dblh{\psline(-0.1,0.05)(0.1,0.05)\psline(-0.1,-0.05)(0.1,-0.05)}
\rput(0.2,4.7){\dblh}\rput(0.2,2.9){\dblh}\rput(0.2,1.1){\dblh}
\rput(5.6,4.7){\dblh}\rput(5.6,2.9){\dblh}\rput(5.6,1.1){\dblh}
\uput[ul](0.2,5.6){A} \uput[ur](5.6,5.6){B} \uput[dr](5.6,0.2){C} \uput[dl](0.2,0.2){D}
\uput[u](2,5.6){I} \uput[u](3.8,5.6){J} \uput[r](5.6,3.8){K} \uput[r](5.6,2){L}
\uput[d](3.8,0.2){M} \uput[d](2,0.2){N} \uput[l](0.2,2){O} \uput[l](0.2,3.8){P}\uput[l](2.9,2.9){S}
\psline[linestyle=dashed,linewidth=1.5pt](0.2,5.6)(5.6,0.2)
\psline[linestyle=dashed,linewidth=1.5pt](5.6,5.6)(0.2,0.2)
\pspolygon(2,5.6)(3.8,5.6)(5.6,3.8)(5.6,2)(3.8,0.2)(2,0.2)(0.2,2)(0.2,3.8)%IJKLMNOP
\end{pspicture}
\end{minipage}
